\section{Application Development}
\label{sec:application_development}

This section documents the end-to-end development of the logistics analytics system,
structured in two phases: a local MVP built on open-source tooling, and a cloud migration
to Google Cloud Platform. Both phases share the same ELT architecture and dbt-derived business
logic; the migration lifts the storage and compute layers without changing the analytical models.

% ─────────────────────────────────────────────────────────────────────────────
\subsection{Local MVP}
\label{sec:local_mvp}

The local pipeline was designed to validate the data model and spatial analysis approach
without cloud dependencies. All components run on a developer machine; the warehouse is a
single DuckDB file at \texttt{db/logistics.duckdb}.

\subsubsection{Pipeline Overview}

\begin{figure}[H]
\centering
\begin{verbatim}
data_generator/
  generate_data.py
        |
        v
data_lake/
  raw_csv/   (customers, products, sellers)
  raw_json/  (orders, delivery_telemetry)
        |
        v
local/load_to_duckdb.py  <-- H3 + grid indices added here
        |
        v
db/logistics.duckdb
        |
        v
local/dbt/               <-- all business logic
  staging/               (views, flatten raw structs)
  intermediate/          (join + cost calculation)
  marts/                 (final tables for BI)
        |
        v
app/app.py               (Streamlit dashboard)
\end{verbatim}
\caption{Local MVP pipeline: data flow from generator to dashboard.}
\label{fig:local_pipeline}
\end{figure}

\subsubsection{Module: Data Generator}
\label{sec:data_generator}

\textbf{File:} \texttt{data\_generator/generate\_data.py}

The data generator produces a synthetic but realistic HCMC logistics dataset using the
\texttt{Faker} library with the \texttt{vi\_VN} locale. It writes two output formats to
\texttt{data\_lake/}:

\begin{itemize}
  \item \textbf{CSV} (\texttt{raw\_csv/}): dimension tables — \texttt{customers.csv},
        \texttt{products.csv}, \texttt{sellers.csv}. Loaded with
        \texttt{read\_csv\_auto} in DuckDB.
  \item \textbf{JSON} (\texttt{raw\_json/}): fact tables — \texttt{orders.json} and
        \texttt{delivery\_telemetry.json}. Loaded as nested \texttt{STRUCT} types with
        \texttt{read\_json\_auto}; fields are flattened only in dbt staging.
\end{itemize}

The generator encodes the primary research variable — spatial blindness — directly into
the data. Eleven HCMC districts are defined with GPS coordinates; four are flagged
\texttt{complex = True} (Quận~3, Quận~5, Quận~10, Gò~Vấp):

\begin{verbatim}
DISTRICTS = [
    {'name': 'Quận 3',  'lat': 10.7801, 'lon': 106.6821, 'complex': True},
    {'name': 'Quận 5',  'lat': 10.7553, 'lon': 106.6639, 'complex': True},
    {'name': 'Quận 10', 'lat': 10.7733, 'lon': 106.6668, 'complex': True},
    {'name': 'Gò Vấp',  'lat': 10.8388, 'lon': 106.6654, 'complex': True},
    ...
]
\end{verbatim}

Deliveries to complex districts receive a 65\% probability of a 15--60 minute delay.
This creates the ground-truth spatial signal that the H3 and grid indexing methods
attempt to recover.

\subsubsection{Module: Local Data Loader}
\label{sec:local_loader}

\textbf{File:} \texttt{local/load\_to\_duckdb.py}

The loader reads all data lake files into DuckDB raw tables and enriches
\texttt{raw\_deliveries} with two spatial index columns computed at ingest time:

\begin{enumerate}
  \item \textbf{H3 hexagonal index} (\texttt{h3\_cell\_9}): computed via a Python UDF
        registered against DuckDB using the \texttt{h3} library. Resolution~9 produces
        cells of approximately 0.1~km², providing sub-block granularity.

        \begin{verbatim}
def py_h3_cell(lat, lon, res=9):
    return h3lib.latlng_to_cell(lat, lon, res)

con.create_function('py_h3_cell', py_h3_cell)
        \end{verbatim}

  \item \textbf{Regular grid index} (\texttt{grid\_cell\_id}): a pure-SQL expression
        that assigns each delivery to a 0.01° × 0.01° cell (approximately 1.1~km²):

        \begin{verbatim}
grid_cell_id = floor(lat/0.01) || '_' || floor(lon/0.01)
        \end{verbatim}
\end{enumerate}

Both indices flow unchanged through staging to the mart layer, where they serve as the
\texttt{GROUP BY} keys in \texttt{mart\_h3} and \texttt{mart\_grid}. The DuckDB community
H3 extension is deliberately not used, so the pipeline runs without external DuckDB plugins.

\subsubsection{Module: dbt Transformation Layer}
\label{sec:dbt}

\textbf{Directory:} \texttt{local/dbt/}

The transformation layer follows a strict ELT pattern: raw data enters the warehouse
untransformed; all business logic lives in dbt SQL. The model DAG is:

\begin{figure}[H]
\centering
\begin{verbatim}
raw_orders       -> stg_orders        \
                 -> stg_order_items    \
raw_deliveries   -> stg_deliveries     -> int_delivery_metrics -> mart_h3
raw_customers    -> stg_customers                                mart_grid
raw_products     -> stg_products                                 mart_kpis
raw_sellers      -> stg_sellers                                  mart_district
                                                                 mart_deliveries
                                                                 mart_spatial_comparison
\end{verbatim}
\caption{dbt model DAG.}
\label{fig:dbt_dag}
\end{figure}

\paragraph{Staging layer.}
Each staging model (e.g., \texttt{stg\_deliveries.sql}, \texttt{stg\_orders.sql}) is
materialized as a \emph{view} and performs a single operation: flattening one raw source.
JSON struct fields are accessed via dot notation:

\begin{verbatim}
telemetry.actual_arrival_at AS actual_arrival_at,
telemetry.spatial_data.destination_lat AS destination_lat
\end{verbatim}

The \texttt{line\_items} array in orders is expanded with a correlated \texttt{UNNEST}.

\paragraph{Intermediate layer.}
\texttt{int\_delivery\_metrics} joins \texttt{stg\_orders} with \texttt{stg\_deliveries}
and computes the core financial signals:

\begin{itemize}
  \item \texttt{delay\_minutes}: \texttt{datediff('minute', estimated, actual)}
  \item \texttt{delay\_cost\_usd}: motorbike \$3/hr, van \$5/hr
  \item \texttt{margin\_gap\_usd}: \texttt{delay\_cost\_usd} $-$ \texttt{total\_freight\_usd};
        a positive value means the delay cost exceeds the freight revenue collected
\end{itemize}

\paragraph{Mart layer.}
Marts are materialized as \emph{tables}. Key outputs:

\begin{itemize}
  \item \texttt{mart\_h3}: one row per H3 cell; aggregates delay, cost, and failure rate;
        flags \texttt{is\_hotspot = true} when \texttt{avg(delay\_minutes) > 15}.
  \item \texttt{mart\_grid}: identical structure, aggregated over 0.01° grid cells.
  \item \texttt{mart\_district}: one row per destination district.
  \item \texttt{mart\_kpis}: single-row summary of fleet-wide metrics.
  \item \texttt{mart\_deliveries}: delivery-level table for map rendering.
  \item \texttt{mart\_spatial\_comparison}: two-row comparison table (H3 vs Grid) with
        five formal spatial metrics documented in Section~\ref{sec:spatial_comparison}.
\end{itemize}

\paragraph{Custom macro.}
A \texttt{haversine\_km} Jinja macro in \texttt{macros/haversine.sql} computes great-circle
distance between two coordinate pairs. It is used in \texttt{mart\_spatial\_comparison} to
calculate the Mean Distance Quantization Error (MDQE) for both indexing methods.

\subsubsection{Module: mart\_spatial\_comparison}
\label{sec:spatial_comparison}

\textbf{File:} \texttt{local/dbt/models/marts/mart\_spatial\_comparison.sql}

This mart produces the primary research output: a side-by-side evaluation of the two
spatial indexing strategies. Each row represents one method; the following metrics are
computed:

\begin{table}[H]
\centering
\begin{tabular}{lp{9cm}}
\hline
\textbf{Metric} & \textbf{Definition} \\
\hline
\texttt{hotspot\_capture\_rate\_pct} & Recall: percentage of all delayed deliveries
  (\texttt{delay > 15 min}) that fall inside flagged hotspot cells. \\
\texttt{hotspot\_precision\_pct} & Precision: percentage of deliveries inside hotspot
  cells that are actually delayed. \\
\texttt{margin\_gap\_capture\_pct} & Percentage of the total margin gap concentrated
  inside flagged cells. \\
\texttt{avg\_within\_cell\_delay\_stddev} & Mean standard deviation of delay within each
  cell; lower values indicate tighter, more homogeneous groupings. \\
\texttt{mdqe\_km} & Mean Distance Quantization Error — average Haversine distance (km)
  from each delivery GPS point to its assigned cell centroid. Measures information lost
  by snapping continuous coordinates to a discrete cell. \\
\texttt{sdc\_km} & Spatial Distortion Coefficient — standard deviation of center-to-neighbour
  distances. H3 = 0.000 by construction (equidistant); Grid $\approx$ 0.228~km at HCMC
  latitudes due to the asymmetry between orthogonal ($\approx$1.11~km) and diagonal
  ($\approx$1.56~km) neighbours. \\
\texttt{cl\_us\_per\_row} & Computational Latency in microseconds per row, benchmarked
  at ingest time (Python UDF for H3 vs pure-SQL \texttt{FLOOR} for Grid). \\
\hline
\end{tabular}
\caption{Spatial comparison metrics computed in \texttt{mart\_spatial\_comparison}.}
\label{tab:spatial_metrics}
\end{table}

\subsubsection{Module: Streamlit Dashboard}
\label{sec:streamlit}

\textbf{File:} \texttt{app/app.py}

The dashboard reads exclusively from mart tables — no aggregations or business logic
are performed in the application layer. It supports both a local DuckDB source and a
BigQuery source, toggled via the \texttt{--source} CLI flag or a sidebar radio button.

The application is structured as four pages:

\begin{itemize}
  \item \textbf{Overview}: five KPI cards (total deliveries, average delay, percentage
        severely delayed, driver cost at risk, margin gap) and two district-level bar charts
        (average delay and margin gap by district). Supports district-level filtering.

  \item \textbf{Delay Map}: a PyDeck \texttt{HeatmapLayer} rendered on a CartoCDN
        dark-matter basemap. Heatmap weight is \texttt{delay\_minutes}, so intensity
        concentrates over the four complex districts. Three summary KPI cards appear below.

  \item \textbf{Districts}: a styled \texttt{st.dataframe} with conditional cell
        colouring (red for severe delay, green for profitable margin), supplemented by a
        failure-rate bar chart and a delay-cost vs margin-gap scatter plot.

  \item \textbf{Hotspots}: H3 cell-level table filtered to \texttt{is\_hotspot = true},
        with three summary KPIs and a pricing recommendation banner.
\end{itemize}

Charts are rendered with Altair and configured with a dark theme matching the dashboard
background. A district multiselect filter in the sidebar applies to the Overview,
Districts, and Hotspots pages.

% ─────────────────────────────────────────────────────────────────────────────
\subsection{Google Cloud Migration}
\label{sec:cloud_migration}

The cloud pipeline reuses the same data lake, ELT pattern, and analytical model structure
as the local MVP. The migration replaces DuckDB with BigQuery and dbt with Dataform,
while GCS serves as the staging area between the data lake and the warehouse.

\subsubsection{Architecture}

\begin{figure}[H]
\centering
\begin{verbatim}
data_lake/
  raw_csv/   raw_json/
        |
        v
cloud/upload_to_gcs.py     <-- JSON -> NDJSON, adds H3 + grid indices
        |
        v
GCS bucket
  raw/csv/   raw/json/
        |
        v
cloud/load_to_bigquery.py  <-- native BQ load jobs
        |
        v
BigQuery: logistics_raw dataset
  (raw_customers, raw_products, raw_sellers,
   raw_orders, raw_deliveries)
        |
        v
cloud/dataform/            <-- SQLX equivalents of dbt models
  staging/   intermediate/   marts/
        |
        v
BigQuery: logistics_marts dataset
        |
        v
app/app.py --source cloud
\end{verbatim}
\caption{Cloud pipeline: GCS staging area between data lake and BigQuery.}
\label{fig:cloud_pipeline}
\end{figure}

\subsubsection{Module: GCS Uploader}
\label{sec:gcs_uploader}

\textbf{File:} \texttt{cloud/upload\_to\_gcs.py}

The uploader performs two roles. First, it converts JSON array files to
newline-delimited JSON (NDJSON), which is required by BigQuery's native load jobs:

\begin{verbatim}
records = json.load(f)          # load full JSON array
for r in records:
    ndjson_lines.append(json.dumps(r))   # one record per line
\end{verbatim}

Second, it enriches each \texttt{delivery\_telemetry} record with the same two spatial
indices computed by the local loader (using the \texttt{h3} Python library and a
\texttt{FLOOR}-based string formula), ensuring parity between the two pipelines:

\begin{verbatim}
record['h3_cell_9']    = h3lib.latlng_to_cell(lat, lon, 9)
record['grid_cell_id'] = f"{math.floor(lat/0.01)}_{math.floor(lon/0.01)}"
\end{verbatim}

CSV dimension files are uploaded directly without transformation.
Required environment variables: \texttt{GCS\_BUCKET}, \texttt{GCP\_PROJECT}.

\subsubsection{Module: BigQuery Loader}
\label{sec:bq_loader}

\textbf{File:} \texttt{cloud/load\_to\_bigquery.py}

Runs BigQuery native load jobs from GCS using the Python client library. No third-party
ingestion tools (Airbyte, Fivetran) are required. Each table is configured with its
schema and the appropriate source format:

\begin{itemize}
  \item CSV tables (\texttt{customers}, \texttt{products}, \texttt{sellers}):
        \texttt{SourceFormat.CSV} with \texttt{skip\_leading\_rows=1}.
  \item JSON tables (\texttt{orders}, \texttt{deliveries}):
        \texttt{SourceFormat.NEWLINE\_DELIMITED\_JSON} using the NDJSON files
        produced by the uploader.
\end{itemize}

All tables are loaded into the \texttt{logistics\_raw} dataset with
\texttt{WriteDisposition.WRITE\_TRUNCATE} (full reload on each run).

\subsubsection{Module: Dataform Transforms}
\label{sec:dataform}

\textbf{Directory:} \texttt{cloud/dataform/}

The Dataform project is the BigQuery-dialect equivalent of the local dbt project.
The model DAG is identical in structure; the key SQL dialect differences are:

\begin{table}[H]
\centering
\begin{tabular}{ll}
\hline
\textbf{DuckDB / dbt} & \textbf{BigQuery / Dataform} \\
\hline
\texttt{datediff('minute', a, b)} & \texttt{TIMESTAMP\_DIFF(b, a, MINUTE)} \\
\texttt{COUNT(*) FILTER (WHERE ...)} & \texttt{COUNTIF(...)} \\
\texttt{MODE()} & \texttt{APPROX\_TOP\_COUNT(col, 1)[OFFSET(0)].value} \\
\texttt{a / b} (may divide by zero) & \texttt{SAFE\_DIVIDE(a, b)} \\
\texttt{\{\{ ref('model') \}\}} & \texttt{\$\{ref('model')\}} \\
\hline
\end{tabular}
\caption{SQL dialect differences between the dbt (DuckDB) and Dataform (BigQuery) layers.}
\label{tab:sql_dialect}
\end{table}

Each SQLX file declares its materialization type in a config block:

\begin{verbatim}
config {
  type: "table",
  schema: "logistics_marts",
  description: "H3 cell-level hotspot aggregation"
}
\end{verbatim}

Staging models use \texttt{type: "view"}; intermediate and mart models use
\texttt{type: "table"}. Source declarations in \texttt{definitions/sources/} register
the raw BigQuery tables as Dataform data sources, replacing dbt's \texttt{sources.yml}.

\subsubsection{Dashboard: Cloud Source}

The Streamlit dashboard connects to BigQuery when launched with \texttt{--source cloud}
or when the sidebar toggle is set to \textit{Cloud --- BigQuery}. The connection is
handled via the \texttt{google-cloud-bigquery} Python client library, authenticated
through either a service account key (\texttt{GOOGLE\_APPLICATION\_CREDENTIALS}) or
Application Default Credentials (\texttt{gcloud auth application-default login}).

Table references switch automatically between the two engines using the \texttt{T()}
helper function:

\begin{verbatim}
def T(name, ds, eng, proj):
    if eng == "duckdb":
        return f"{ds}.{name}"
    return f"`{proj}.{ds}.{name}`"
\end{verbatim}

All queries, chart logic, and filtering behaviour are identical between the local and
cloud sources; only the connection backend and table reference format differ.

% ─────────────────────────────────────────────────────────────────────────────
\subsection{Design Decisions and Trade-offs}
\label{sec:design_decisions}

\paragraph{Spatial indices computed at load time, not in dbt.}
Both the H3 UDF and the grid formula are applied in the loader (or uploader), not in
dbt SQL. This ensures the indices are present in the raw table and available to all
downstream models without repeating the computation. It also keeps the dbt layer
purely declarative: no Python models or external function calls are required.

\paragraph{ELT over ETL.}
Raw data enters the warehouse unchanged. Transformation happens in SQL inside the
warehouse, which makes it auditable, version-controlled, and independently testable.
The data lake files serve as the immutable source of truth.

\paragraph{Hotspot threshold hardcoded at 15 minutes.}
The \texttt{is\_hotspot} flag is set where \texttt{avg(delay\_minutes) > 15}.
This threshold appears in both \texttt{mart\_h3} and \texttt{mart\_grid} and must be
changed in both files if adjusted. The value was chosen to match the 65\% delay
injection rate used in the generator (15--60 minute injected delays against a baseline
that produces near-zero delay for non-complex districts).

\paragraph{No dbt for cloud; Dataform instead.}
Dataform is the native transformation framework for BigQuery and integrates directly
with the Google Cloud ecosystem (IAM, Workflows, Cloud Scheduler). Using dbt against
BigQuery would have required an additional adapter and separate orchestration; Dataform
eliminates this at the cost of a BigQuery-specific syntax.
